\documentclass[11pt,a4paper]{article}

\usepackage[utf8]{inputenc}
\usepackage[T1]{fontenc}
\usepackage[spanish,es-noquoting]{babel}
\usepackage[margin=2.5cm]{geometry}
\usepackage{amsmath}
\usepackage{booktabs}
\usepackage{hyperref}

\hypersetup{colorlinks=true, linkcolor=black, urlcolor=blue}

\title{Cómo funciona \texttt{receptor.py}\\
       \large Documentación del módulo de balance nodal del receptor}
\author{Modelo del receptor solar PPHE}
\date{19 de septiembre de 2026}

\begin{document}
\maketitle
\tableofcontents
\bigskip

\section{Qué hace el módulo y dónde encaja}

\texttt{receptor.py} es el módulo que \emph{une} las dos mitades del modelo. No
contiene correlaciones propias --- por eso el \texttt{formulas.py} de esta carpeta
sigue vacío ---: lo que hace es tomar un mapa de flujo incidente, un fluido con
propiedades variables y la cadena termohidráulica del canal de una
\emph{pillow plate}, y resolver con las tres cosas el balance de energía de la
placa, porción a porción.

Su entrada es un objeto \texttt{MapaGaussiano} más un punto de trabajo (gasto,
temperatura de entrada, presión, malla, geometría de panel). Su salida es un
objeto \texttt{Resultado} con dos matrices de temperatura --- la del fluido y la de
la pared --- y el calor total absorbido. El resultado que pedía el encargo del
modelo es el perfil $T(x)$ en $y = L$.

\subsection{Dependencias}

\begin{center}
\begin{tabular}{lll}
  \toprule
  Módulo & Se importa como & Qué aporta \\
  \midrule
  \texttt{modelo\_receptor/mapa.py}          & \texttt{mp}   & $q''(x,y)$ y su reparto nodal \\
  \texttt{modelo\_receptor/props\_fluido.py} & \texttt{pf}   & $\rho$, $c_p$, $k$, $\mu$, $h(T)$ del aire \\
  \texttt{modelo\_termohidraulico/formulas.py}   & \texttt{form} & $Re$, $Nu_I$ (20), $d_{e,I}$ (14) \\
  \texttt{modelo\_termohidraulico/asignacion.py} & \texttt{asig} & constantes $n_1\ldots n_5$ \\
  \texttt{modelo\_termohidraulico/geometria.py}  & \texttt{geom} & PPHE1, PPHE2, PPHE3 \\
  \bottomrule
\end{tabular}
\end{center}

\subsection{El rodeo de \texttt{importlib}}

Los tres módulos del otro paquete no se cargan con un \texttt{import} normal,
sino con la función privada \texttt{\_importar(nombre, ruta)}, que construye la
especificación con \texttt{spec\_from\_file\_location} y la ejecuta a mano:

\begin{verbatim}
    spec   = spec_from_file_location(nombre, ruta)
    modulo = module_from_spec(spec)
    spec.loader.exec_module(modulo)
\end{verbatim}

El motivo es un choque de nombres real: las dos carpetas tienen un
\texttt{formulas.py}. Con un \texttt{import formulas} normal, Python cogería el
de la carpeta desde la que se ejecute el script, que unas veces sería el bueno y
otras no, y el fallo sería silencioso --- funciones con el mismo nombre y
distinto contenido ---. Cargando por ruta absoluta se elige sin ambigüedad, y
además no se toca \texttt{sys.path}, de modo que el módulo no deja efectos
secundarios en quien lo importe. Los nombres internos que se les da
(\texttt{th\_formulas}, \texttt{th\_asignacion}, \texttt{th\_geometria}) llevan
prefijo por la misma razón.

La ruta se calcula con \texttt{Path(\_\_file\_\_).resolve().parent.parent}, es
decir, relativa al propio fichero y no al directorio de trabajo: el módulo
funciona se ejecute desde donde se ejecute.

\section{Sistema de coordenadas y discretización}

El origen está en la esquina inferior izquierda de la placa vista de frente,
$x$ hacia la derecha en la anchura y $y$ hacia arriba en la altura, que es la
dirección en que corre el fluido:
\[
  x \in [0, W], \qquad y \in [0, L].
\]

La placa se parte en $N$ filas de altura $\Delta y = L/N$ y $M$ columnas de
anchura $\Delta x = W/M$. Cada porción tiene área
\[
  A = \frac{W}{M}\cdot\frac{L}{N},
\]
y cada columna recibe una fracción $G_{\mathrm{col}} = G/M$ del gasto total.

\paragraph{Cada columna es un canal.} El aire entra por abajo a
$T_{\mathrm{ent}}$ uniforme, atraviesa las $N$ porciones de su columna y sale por
$y = L$. La salida de una porción es la entrada de la de encima. Las columnas no
se hablan entre sí: son adiabáticas unas respecto a otras.

\paragraph{Indexado.} Las matrices de resultado se recorren
\texttt{T[j][i]}, con $j$ creciendo en $y$ (de abajo arriba) e $i$ creciendo en
$x$. Es el mismo convenio que devuelve \texttt{mapa.mapa\_nodal}, de modo que
\texttt{flujo[j][i]} y \texttt{T\_fluido[j][i]} se refieren siempre a la misma
porción. \texttt{T\_fluido[-1]} es, por tanto, la fila de arriba, es decir, el
perfil de salida.

\section{Constantes del entorno y del material}

\subsection{El sumidero radiante}

\begin{center}
\begin{tabular}{llr}
  \toprule
  Constante & Significado & Valor \\
  \midrule
  \texttt{SIGMA}        & Stefan--Boltzmann          & $5{,}670374419\cdot10^{-8}$~W/(m$^2$K$^4$) \\
  \texttt{T\_AMB}       & Temperatura ambiente       & 298,15~K \\
  \texttt{T\_ENTORNO}   & Entorno próximo, 30~$^\circ$C & 303,15~K \\
  \texttt{PESO\_ENTORNO}& Peso del entorno próximo   & 0,895 \\
  \texttt{PESO\_CIELO}  & Peso del cielo             & 0,955 \\
  \bottomrule
\end{tabular}
\end{center}

La temperatura del sumidero no es la del ambiente ni la del cielo por separado,
sino la media en $T^4$ de las dos, ponderada con los factores de la ecuación (6)
del artículo de M.R.\ Rodríguez:
\[
  T_{\mathrm{cielo}} = \left(
    \frac{w_e\,T_{\mathrm{entorno}}^4 + w_c\,T_{\mathrm{amb}}^4}{w_e + w_c}
  \right)^{1/4} = 300{,}60~\mathrm{K} = 27{,}45~^\circ\mathrm{C}.
\]
La media se hace en la cuarta potencia, y no en la temperatura, porque lo que se
suma físicamente son emitancias, que van con $T^4$; promediar temperaturas daría
otro número.

Es importante notar que este valor se calcula \emph{una vez} al importar el
módulo y queda fijo. No es un parámetro del constructor, a diferencia de
$h_{\mathrm{ext}}$, $\varepsilon$ y $\alpha$: para cambiarlo hay que editar el
módulo.

\subsection{Los tres parámetros de la cara expuesta}

\begin{center}
\begin{tabular}{llrl}
  \toprule
  Constante & Significado & Valor & ¿Configurable? \\
  \midrule
  \texttt{H\_EXT}        & Convección exterior & 20~W/(m$^2$K) & sí, campo \texttt{h\_ext} \\
  \texttt{EPSILON}       & Emisividad          & 0,85          & sí, campo \texttt{eps} \\
  \texttt{ABSORTIVIDAD}  & Absortividad        & 1,00          & sí, campo \texttt{alfa} \\
  \bottomrule
\end{tabular}
\end{center}

$h_{\mathrm{ext}}$ mezcla convección natural y forzada por viento sobre una pared
vertical caliente. La bibliografía de receptores de torre maneja de 10~W/(m$^2$K)
en calma a 30 con viento, y 20 es el valor intermedio habitual. Es de los
parámetros más inciertos del modelo y, aun así, de los menos influyentes, porque
a estas temperaturas la radiación domina: en el nodo más caliente del caso de
referencia se lleva 2,84 veces más calor que la convección, y la proporción
crece con la temperatura (tabla~\ref{tab:radconv}).

\begin{table}[h]
\centering
\begin{tabular}{rrrr}
  \toprule
  $T_w$ [K] & Convección [W/m$^2$] & Radiación [W/m$^2$] & Cociente \\
  \midrule
   800 & 10\,037 & 19\,348 & 1,93 \\
  1000 & 14\,037 & 47\,805 & 3,41 \\
  1150 & 17\,037 & 83\,905 & 4,92 \\
  \bottomrule
\end{tabular}
\caption{Reparto de las pérdidas de la cara expuesta con $h_{\mathrm{ext}} = 20$
y $\varepsilon = 0{,}85$. La radiación es la que manda, y es lo que hace que la
incertidumbre de $h_{\mathrm{ext}}$ importe poco.}
\label{tab:radconv}
\end{table}

$\alpha = 1$ deja el modelo en el balance exacto que pedía el encargo, con todo
el flujo del mapa entrando en la pared. Con una pintura selectiva tipo Pyromark
sería 0,95, y basta cambiar el número: el efecto sobre el caso de referencia es
bajar el rendimiento de 73,28 a 69,49~\% y la pared de 659 a 642~$^\circ$C.

\subsection{Conductividad del acero}
\label{sec:kacero}

\[
  K_{\mathrm{acero}} = 20~\mathrm{W/(m\,K)}
\]

Es una \emph{constante}, no una función de la temperatura, y la decisión merece
explicarse, porque el AISI~321 sí depende de ella: conduce unos 16~W/(m\,K) a
100~$^\circ$C y 22 a 600, y los austeníticos, al revés que la mayoría de los
metales, \emph{mejoran} al calentarse.

Lo que permite congelarla es el peso que tiene la chapa en el circuito térmico.
En el nodo de pared más caliente del caso de referencia:

\begin{center}
\begin{tabular}{lrr}
  \toprule
  Resistencia & Valor [m$^2$K/W] & Fracción \\
  \midrule
  Chapa, $e/K_{\mathrm{acero}}$   & $4{,}000\cdot10^{-5}$ & 2,92~\% \\
  Película interna, $1/h_{int}$   & $1{,}328\cdot10^{-3}$ & 97,08~\% \\
  \midrule
  Total                           & $1{,}368\cdot10^{-3}$ & 100~\% \\
  \bottomrule
\end{tabular}
\end{center}

La chapa es el 2,9~\% de la resistencia: de los 184,4~K de salto entre pared y
fluido, se lleva 5,4.

\paragraph{De dónde sale el 20.} Sobre todos los casos de trabajo ensayados
--- picos de 200 y 400~kW/m$^2$, entradas de 100 a 500~$^\circ$C, gastos de 0,3 a
1,5~kg/s y los tres paneles --- la temperatura media de la chapa, $(T_w+T_f)/2$,
se mueve entre 396 y 1078~K, es decir, $k$ entre 16,2 y 24,8~W/(m\,K), con media
20,5. Se toma 20, que corresponde a unos 700~K de chapa, el centro del rango. Un
20~\% de error en $k$ mueve la resistencia total un 0,5~\%.

\paragraph{Cuánto cuesta la simplificación.} Medido contra la versión con
$k(T)$, sobre siete puntos de trabajo, la desviación nodal máxima es de 1,2~K, y
se da en el caso más caliente, el de 400~kW/m$^2$ de pico. En el de referencia
es de 0,34~K en la pared máxima, 0,02~K en la salida media y 0,009 puntos en el
rendimiento. Para ponerlo en contexto: pasar $h_{\mathrm{ext}}$ de 20 a
10~W/(m$^2$K), dentro de su propio rango de incertidumbre, mueve la pared 11~K.
La dependencia de $k$ con $T$ estaba muy por debajo del ruido del modelo.

\paragraph{Y además desacopla la bisección.} Ésta es la razón de fondo. Con
$k = k(T)$ la resistencia dependía de $T_w$ --- la temperatura media de la chapa
es $(T_w + T_f)/2$, y $T_w$ es justo la incógnita ---, así que había que
reevaluarla en cada iteración de la bisección interna. Con $k$ constante la
resistencia se calcula una vez por porción y sale del bucle
(sección~\ref{sec:coste}).

El límite de servicio continuo del AISI~321 está en torno a 1150~K
(877~$^\circ$C). El módulo no lo impone: lo avisa en \texttt{resumen()}.

\section{Las dos funciones de módulo}

\subsection{\texttt{perdidas(T\_w, A, h\_ext, eps)}}

\[
  q_{\mathrm{perd}} = \Bigl[\,h_{\mathrm{ext}}\,(T_w - T_{\mathrm{amb}})
    + \varepsilon\,\sigma\,(T_w^4 - T_{\mathrm{cielo}}^4)\,\Bigr]\,A
  \qquad [\mathrm{W}]
\]

Está \emph{fuera} de la clase a propósito. Se llama desde el interior de la
bisección, del orden de mil veces por porción, y sacarla del objeto permite
pasarle $h_{\mathrm{ext}}$ y $\varepsilon$ como argumentos locales en vez de
resolver dos accesos a atributos en cada llamada.

\subsection{\texttt{\_biseccion(f, a, b, tol=1e-6, iteraciones=80)}}

Bisección clásica, con dos detalles que merecen comentario.

\paragraph{Exige que el intervalo encierre la raíz.} Evalúa $f(a)$ y $f(b)$ y
lanza \texttt{ValueError} si tienen el mismo signo, en vez de devolver un
resultado sin sentido. Quien llama es responsable de acotar bien, y las dos
llamadas del módulo lo hacen con argumentos físicos (sección~\ref{sec:porcion}).

\paragraph{Una sola evaluación por iteración.} El bucle guarda \texttt{fa} y lo
actualiza solo cuando mueve el extremo izquierdo:

\begin{verbatim}
    m  = 0.5 * (a + b)
    fm = f(m)                 # la unica evaluacion de la iteracion
    if fa * fm <= 0:
        b = m                 # fb no hace falta: no se vuelve a usar
    else:
        a, fa = m, fm
\end{verbatim}

Esto no es cosmético. Una versión que reevaluase $f$ para decidir el signo
duplicaría el coste del módulo entero, porque $f$ es aquí una bisección anidada
completa.

\paragraph{Criterio de parada.} Sale cuando $b - a < 10^{-6}$~K, o al agotar las
80 iteraciones. Con los intervalos que usa el módulo, de 300 a 1100~K de
anchura, hacen falta $\log_2(10^9) \approx 30$ iteraciones, así que el tope de 80
nunca se alcanza: es una red de seguridad, no el criterio real. La medida sobre
el caso de referencia da 30,1 iteraciones de media.

\section{La clase \texttt{Receptor}}

Es un \texttt{dataclass} congelado (\texttt{frozen=True}): una vez construido no
se puede modificar, lo que garantiza que los campos derivados que calcula el
constructor siguen siendo coherentes con los de entrada.

\subsection{Campos}

\begin{center}
\begin{tabular}{llll}
  \toprule
  Campo & Tipo & Por defecto & Significado \\
  \midrule
  \texttt{mapa}   & \texttt{MapaGaussiano} & --- & mapa de flujo incidente \\
  \texttt{G}      & float & --- & gasto másico total [kg/s] \\
  \texttt{T\_ent} & float & --- & entrada, uniforme [K] \\
  \texttt{p}      & float & $10\cdot10^5$ & presión del aire [Pa] \\
  \texttt{N}      & int   & 20 & porciones en altura \\
  \texttt{M}      & int   & 10 & porciones en anchura \\
  \texttt{panel}  & obj   & \texttt{PPHE1} & geometría de la pillow plate \\
  \texttt{fluido} & obj   & \texttt{AIRE}  & fluido de trabajo \\
  \texttt{h\_ext} & float & 20,0 & convección exterior \\
  \texttt{eps}    & float & 0,85 & emisividad \\
  \texttt{alfa}   & float & 1,00 & absortividad \\
  \midrule
  \texttt{n}       & obj   & \multicolumn{2}{l}{\emph{derivado}: constantes $n_1\ldots n_5$} \\
  \texttt{seccion} & float & \multicolumn{2}{l}{\emph{derivado}: sección de paso [m$^2$]} \\
  \texttt{de}      & float & \multicolumn{2}{l}{\emph{derivado}: diámetro equivalente [m]} \\
  \bottomrule
\end{tabular}
\end{center}

\subsection{\texttt{\_\_post\_init\_\_}: lo que solo depende de la geometría}

Los tres últimos campos llevan \texttt{field(init=False)}: no se pasan al
construir, los calcula \texttt{\_\_post\_init\_\_} una sola vez. Como la clase
está congelada, hay que escribirlos con \texttt{object.\_\_setattr\_\_}, y de ahí
la \texttt{lambda fijar} que hace ese rodeo legible.

\begin{align}
  n_1 \ldots n_5 &= \texttt{asignacion\_de\_constantes}(s_T,\, 2s_L,\, d_{sp},\, b_i) \\
  A_{\mathrm{paso}} &= \frac{b_i}{\sqrt{2}}\,W && \text{ec.\ (16)} \\
  d_{e,I} &= 2\,\frac{b_i}{\sqrt{2}} && \text{ec.\ (14)}
\end{align}

Sobre la sección de paso hay que hacer una advertencia: se calcula con la
anchura \emph{de la placa}, \texttt{self.mapa.W}, y no con la anchura útil del
panel $w_{pp} - 2w_e$ que usa \texttt{formulas.f\_chI}. Es decir, se supone que
el canal ocupa toda la anchura del receptor y se desprecian las soldaduras de
borde. Es coherente con el resto del modelo, que trata la placa como una única
pillow plate de $L\times W$, pero no es lo mismo que apilar paneles de
$w_{pp} = 300$~mm.

Nótese también que la asignación de constantes se hace en el constructor, de modo
que una geometría fuera de las correlaciones publicadas falla al crear el objeto
y no a mitad del cálculo.

\subsection{\texttt{h\_interno(T)}: la cadena termohidráulica}

Devuelve la terna $(h, u, Re)$, siguiendo la misma cadena que
\texttt{modelo\_f3.py} del otro paquete:
\begin{align}
  u  &= \frac{G}{\rho(T,p)\,A_{\mathrm{paso}}}  && \text{ec.\ (3)} \\
  Re &= \frac{\rho\,u\,d_{e,I}}{\mu}            && \\
  Nu &= n_3\,Re^{\,n_4}\,Pr^{\,n_5}             && \text{ec.\ (20)} \\
  h  &= \frac{Nu\,k}{d_{e,I}}                   &&
\end{align}

La diferencia con \texttt{modelo\_f3} es que aquí las propiedades \emph{no} son
constantes: se piden a \texttt{props\_fluido} a la temperatura que se le pase, y
el módulo la llama en cada porción con la temperatura media del fluido de esa
porción. En el caso de referencia, $h$ pasa de 700 a 772~W/(m$^2$K) de abajo
arriba de la columna central.

\paragraph{Por qué $h$ no depende de $M$.} El gasto de una columna,
$G_{\mathrm{col}} = G/M$, y su sección de paso, $A_{\mathrm{paso}}/M$, son ambos
proporcionales a $\Delta x$, así que el cociente se cancela y la velocidad es la
misma con cualquier $M$. Por eso \texttt{h\_interno} usa el gasto \emph{total} y
la sección \emph{total}, sin dividir por $M$: es el mismo número. La consecuencia
práctica es que afinar la malla cambia la resolución del mapa, no la física del
canal.

\section{El balance de una porción: \texttt{\_porcion}}
\label{sec:porcion}

Es el corazón del módulo. Resuelve una sola porción y devuelve
$(T_{\mathrm{sal}}, T_w, q_{\mathrm{fluido}})$.

\subsection{Las tres ecuaciones}

Sobre la cara exterior incide $q_{\mathrm{inc}} = q''\!\cdot\!A$, de lo cual se
absorbe $q_{\mathrm{abs}} = \alpha\,q_{\mathrm{inc}}$. La pared, a $T_w$, pierde
al ambiente y al cielo, y lo que queda atraviesa la chapa y pasa al fluido:
\begin{align}
  q_{\mathrm{abs}} &= q_{\mathrm{perd}}(T_w) + q_{\mathrm{fluido}}, \label{eq:bal}\\[2pt]
  q_{\mathrm{fluido}} &= \frac{A\,(T_w - T_f)}
       {\dfrac{e}{K_{\mathrm{acero}}} + \dfrac{1}{h_{\mathrm{int}}(T_f)}},
  \qquad T_f = \tfrac{1}{2}\left(T_{\mathrm{ent}} + T_{\mathrm{sal}}\right), \label{eq:cond}\\[2pt]
  q_{\mathrm{fluido}} &= G_{\mathrm{col}}\left[h(T_{\mathrm{sal}}) - h(T_{\mathrm{ent}})\right].
  \label{eq:ent}
\end{align}

Dos incógnitas, $T_w$ y $T_{\mathrm{sal}}$, y dos no linealidades: el $T_w^4$ de
la radiación y la dependencia de las propiedades del fluido con $T_f$. La
conductividad del acero es constante (sección~\ref{sec:kacero}), de modo que el
denominador de \eqref{eq:cond} no depende de $T_w$ y se evalúa una sola vez por
porción, fuera de la bisección interna.

\paragraph{Por qué entalpías y no $c_p\,\Delta T$.} La ecuación~\eqref{eq:ent}
usa la integral analítica del polinomio de $c_p$ que trae
\texttt{props\_fluido}. Entre 300 y 1200~K el $c_p$ del aire sube un 17~\%, así
que cerrar el balance con el $c_p$ de la entrada se dejaría un 7,4~\% del calor
absorbido por el camino. Con saltos de 200 a 400~K por columna, eso no es un
detalle.

\paragraph{La temperatura de película.} El fluido de una porción se caracteriza
por su temperatura \emph{media} $T_f = (T_{\mathrm{ent}} + T_{\mathrm{sal}})/2$,
y la chapa por la media entre pared y fluido, $(T_w + T_f)/2$. Son las dos
medias aritméticas habituales, y son las que hacen que las ecuaciones sean
implícitas: $T_f$ depende de $T_{\mathrm{sal}}$, que es lo que se busca.

\subsection{La estructura: bisección anidada}

\begin{verbatim}
    _porcion(q_inc, T_ent, G_col, A)
      |
      +-- biseccion EXTERNA sobre T_sal, en [T_min_fluido, T_max]
      |     f = desequilibrio(T_sal) = q_fluido(T_sal) - G_col*[h(T_sal)-h(T_ent)]
      |     |
      |     +-- cerrar(T_sal):
      |           T_f = (T_ent + T_sal)/2
      |           h_int = h_interno(T_f)
      |           R     = e/K_ACERO + 1/h_int     # constante dentro del bucle
      |           |
      |           +-- biseccion INTERNA sobre T_w, en [T_cielo, T_w_max]
      |                 f = balance_pared(T_w) = q_abs - perdidas(T_w) - conduccion(T_w)
      |
      +-- cerrar(T_sal_final) una ultima vez, para devolver T_w y q
\end{verbatim}

Para cada $T_{\mathrm{sal}}$ de prueba se fija $T_f$, con él se calcula
$h_{\mathrm{int}}$, se despeja $T_w$ de la ecuación~\eqref{eq:bal} con
\eqref{eq:cond} sustituida, y con ese $T_w$ se obtiene el calor al fluido. Ese
calor se compara con el salto de entalpía \eqref{eq:ent}, y el desequilibrio es
la función que anula la bisección externa.

\subsection{El intervalo de la bisección interna, sobre $T_w$}

\[
  T_w \in \Bigl[\;T_{\mathrm{cielo}},\;\;
  \max\Bigl(T_f,\ \bigl(\tfrac{q_{\mathrm{abs}}}{A\,\varepsilon\,\sigma} + T_{\mathrm{cielo}}^4\bigr)^{1/4}\Bigr) + 1\;\Bigr]
\]

Los dos extremos tienen lectura física:

\begin{itemize}
  \item \textbf{Extremo inferior.} A $T_w = T_{\mathrm{cielo}}$ la radiación neta
        es nula y la pared, más fría que el fluido, recibe calor por los dos
        lados: el balance es positivo.
  \item \textbf{Extremo superior.} Es la temperatura a la que \emph{solo} la
        radiación ya evacúa todo el flujo absorbido. Como además hay convección
        y conducción al fluido, el balance allí es necesariamente negativo. El
        \texttt{max} con $T_f$ cubre el caso de flujo casi nulo, en que la
        expresión de la radiación daría un valor por debajo del fluido, y el
        $+1$ evita el intervalo degenerado cuando los dos extremos coinciden.
\end{itemize}

En el caso de referencia, el intervalo de la celda más caliente es
$[300{,}6;\ 1422{,}3]$~K, 1121,7~K de anchura, que con tolerancia $10^{-6}$~K
pide 30,1 iteraciones.

\subsection{El intervalo de la bisección externa, sobre $T_{\mathrm{sal}}$}

\[
  T_{\mathrm{sal}} \in \Bigl[\;T_{\min}^{\mathrm{fluido}},\;\;
  \min\Bigl(T_{\mathrm{ent}} + \frac{q_{\mathrm{abs}}}{G_{\mathrm{col}}\,c_p(T_{\mathrm{ent}})},\ T_{\max}^{\mathrm{fluido}}\Bigr)\;\Bigr]
\]

El extremo superior supone que \emph{todo} el calor absorbido va al fluido, sin
ninguna pérdida, y lo evalúa con el $c_p$ de la entrada, que es el menor del
tramo. Como el $c_p$ del aire crece con la temperatura, la temperatura real
queda necesariamente por debajo: la cota es rigurosa. El \texttt{min} con
$T_{\max}^{\mathrm{fluido}} = 1600$~K impide pedir a \texttt{props\_fluido} una
temperatura fuera de su rango de ajuste.

El extremo inferior son los 250~K del ajuste, no $T_{\mathrm{ent}}$. Es
deliberado: permite que el fluido se \emph{enfríe} en una porción, cosa que pasa
de verdad cuando el flujo incidente es pequeño y el aire entra caliente.

\paragraph{La comprobación previa.} Antes de bisecar, el código evalúa
\texttt{desequilibrio(T\_max)} y, si sale positivo, lanza \texttt{ValueError}
con un mensaje legible en vez de dejar que \texttt{\_biseccion} falle con su
error genérico. Un desequilibrio positivo en el extremo superior significa que
ni siquiera a 1600~K el fluido absorbe lo que le llega: el punto de trabajo no
existe, y el consejo del mensaje --- subir el gasto o bajar el flujo --- es el
correcto.

\subsection{Por qué bisección y no Newton}

Newton convergería en tres o cuatro iteraciones en vez de treinta, pero:

\begin{itemize}
  \item habría que derivar el balance, y la derivada incluye la de
        $k_{\mathrm{acero}}$, la de $h_{\mathrm{int}}$ con la temperatura y la
        del polinomio de entalpía;
  \item puede divergir, y aquí se ejecuta dentro de un bucle de cientos de
        porciones sin nadie mirando;
  \item el coste no es el cuello de botella con mallas de este tamaño.
\end{itemize}

La bisección, en cambio, converge siempre que el intervalo encierre la raíz, y
los dos intervalos se han acotado con argumentos físicos que lo garantizan
(salvo el caso de la sección~\ref{sec:fallo}).

\subsection{Coste: recuento de evaluaciones}
\label{sec:coste}

Medido sobre el caso de referencia, malla $10\times20$, instrumentando
\texttt{\_biseccion}:

\begin{center}
\begin{tabular}{lr}
  \toprule
  Porciones                                  & 200 \\
  Bisecciones totales                        & 6\,820 \\
  Bisecciones por porción                    & 34,1 \\
  Iteraciones por bisección                  & 30,1 \\
  Evaluaciones de \texttt{balance\_pared}    & 218\,852 \\
  Evaluaciones por porción                   & 1\,094 \\
  \bottomrule
\end{tabular}
\end{center}

Las 34 bisecciones internas por porción se reparten así: 32 de la bisección
externa (30 iteraciones más las 2 evaluaciones del chequeo de signo), 1 de la
comprobación previa de \texttt{T\_max} y 1 de la llamada final a
\texttt{cerrar}. Cada una de ellas cuesta a su vez unas 32 evaluaciones del
balance de pared, de donde salen los $\sim$1100 por porción.

\paragraph{Optimizaciones presentes.} \texttt{\_porcion} saca a variables
locales todo lo que necesita del objeto --- \texttt{fluido}, \texttt{e},
\texttt{h\_ext}, \texttt{eps}, \texttt{h\_interno} --- antes de definir las
funciones internas. Con 200\,000 evaluaciones por placa, cada acceso a atributo
que se evita se nota.

En la misma línea, la resistencia térmica se calcula una vez por llamada a
\texttt{cerrar} y no dentro de \texttt{balance\_pared}, cosa que solo es posible
porque $K_{\mathrm{acero}}$ es constante. El recuento de evaluaciones no cambia
por ello --- siguen siendo 218\,852 ---: lo que baja es el coste de cada una, y con
él un 18~\% el tiempo total.

\section{El recorrido de la placa: \texttt{resolver}}

\begin{verbatim}
    A     = (W/M) * (L/N)
    G_col = G / M
    flujo = mapa.mapa_nodal(M, N)      # [W/m2] medio de cada celda

    para cada columna i:
        T = T_ent                      # entrada uniforme
        para cada fila j, de abajo arriba:
            T, T_w, q = porcion(flujo[j][i] * A, T, G_col, A)
            T_fluido[j][i] = T         # la salida de esta es la entrada de la de encima
            T_pared[j][i]  = T_w
            q_util += q
\end{verbatim}

El bucle exterior recorre columnas y el interior filas, en ese orden, porque el
acoplamiento es \emph{vertical}: la temperatura se arrastra hacia arriba dentro
de una columna, y no hay nada que arrastrar entre columnas. Es un barrido
explícito, sin iteración global: cada porción se resuelve una vez y su resultado
es definitivo.

El flujo nodal viene de \texttt{mapa.mapa\_nodal(M, N)}, que \emph{integra} la
gaussiana dentro de cada celda con la función error en vez de muestrearla en su
centro. Eso hace que $\sum q''_{ij}\,A$ reproduzca exactamente la potencia total
sobre la placa, y es lo que permite que el balance global cierre con mallas
bastas.

\section{La clase \texttt{Resultado}}

Contenedor congelado con las dos matrices, el calor útil y una referencia al
receptor que lo produjo.

\begin{center}
\begin{tabular}{lll}
  \toprule
  Miembro & Tipo & Qué da \\
  \midrule
  \texttt{T\_fluido[j][i]} & matriz & salida de cada porción [K] \\
  \texttt{T\_pared[j][i]}  & matriz & temperatura de pared [K] \\
  \texttt{q\_util}         & float  & calor total al fluido [W] \\
  \texttt{perfil\_salida}  & prop.  & \texttt{T\_fluido[-1]}, el perfil en $y=L$ \\
  \texttt{x\_centros}      & prop.  & $x$ del centro de cada columna [m] \\
  \texttt{T\_pared\_max}   & prop.  & máximo de la matriz de pared [K] \\
  \texttt{rendimiento}     & prop.  & \texttt{q\_util} / potencia incidente \\
  \texttt{resumen()}       & método & imprime el balance y el perfil \\
  \texttt{dibujar()}       & método & campo de fluido y perfil, con matplotlib \\
  \bottomrule
\end{tabular}
\end{center}

\texttt{resumen()} recalcula $h$, $u$ y $Re$ a la temperatura media entre la
entrada y la media del perfil de salida, solo para informar, y marca con un aviso
explícito si la pared pasa de 1150~K.

\texttt{dibujar()} saca dos paneles: el campo de temperatura del fluido con
\texttt{imshow}, escala \texttt{hot}, \texttt{origin="lower"} y
\texttt{interpolation="nearest"} --- sin interpolar, para que la discretización se
vea ---, y el perfil de salida. Importa matplotlib dentro de la función y avisa sin
romper si no está instalado, de modo que el módulo corre en un entorno sin
gráficos.

\section{Comportamiento numérico}

\subsection{Convergencia de malla}

\begin{table}[h]
\centering
\begin{tabular}{lrrrrr}
  \toprule
  $M\times N$ & $Q_{\mathrm{útil}}$ [kW] & $\eta$ [\%] & $\bar{T}_{\mathrm{sal}}$ [$^\circ$C]
    & $T_{\mathrm{sal}}^{\max}$ [$^\circ$C] & $T_w^{\max}$ [$^\circ$C] \\
  \midrule
  $5\times10$  & 115,3038 & 73,344 & 479,944 & 545,57 & 658,09 \\
  $10\times20$ & 115,1887 & 73,270 & 479,728 & 545,51 & 659,35 \\
  $20\times40$ & 115,1596 & 73,252 & 479,674 & 548,11 & 663,18 \\
  $40\times80$ & 115,1523 & 73,247 & 479,660 & 548,77 & 664,25 \\
  \bottomrule
\end{tabular}
\caption{Independencia de malla del caso de referencia.}
\label{tab:malla}
\end{table}

Las magnitudes integrales convergen muy deprisa: de $5\times10$ a $40\times80$ el
rendimiento se mueve 0,10 puntos y la temperatura media de salida, 0,28~K. Los
\emph{picos} convergen más despacio --- la pared máxima sube 6,1~K --- porque una
malla basta promedia la cresta de la campana dentro de una celda y rebaja el
punto caliente. Con $10\times20$ ya se está a 5~K del valor convergido de la
pared, y esa es la razón de que sea la malla por defecto.

\subsection{Coste computacional}

\begin{center}
\begin{tabular}{lrr}
  \toprule
  Malla & Porciones & Tiempo [s] \\
  \midrule
  $5\times10$  &    50 & 0,03 \\
  $10\times20$ &   200 & 0,13 \\
  $20\times40$ &   800 & 0,51 \\
  $40\times80$ & 3\,200 & 2,07 \\
  \bottomrule
\end{tabular}
\end{center}

El coste es lineal en el número de porciones, como cabe esperar de un barrido
explícito en el que cada porción cuesta lo mismo: unas 1100 evaluaciones del
balance de pared.

\subsection{Sensibilidad a los parámetros inciertos}

\begin{center}
\begin{tabular}{lrrr}
  \toprule
  & $\eta$ [\%] & $T_w^{\max}$ [$^\circ$C] & $\bar{T}_{\mathrm{sal}}$ [$^\circ$C] \\
  \midrule
  $h_{\mathrm{ext}} = 10$              & 76,83 & 670,8 & 488,3 \\
  $h_{\mathrm{ext}} = 20$ \emph{(base)}& 73,27 & 659,3 & 479,7 \\
  $h_{\mathrm{ext}} = 30$              & 69,83 & 648,1 & 471,4 \\
  \midrule
  $\varepsilon = 0{,}70$               & 75,73 & 670,1 & 485,6 \\
  $\varepsilon = 0{,}85$ \emph{(base)} & 73,27 & 659,3 & 479,7 \\
  $\varepsilon = 1{,}00$               & 70,97 & 649,3 & 474,2 \\
  \midrule
  $\alpha = 0{,}95$ (Pyromark)         & 69,48 & 642,7 & 470,6 \\
  $\alpha = 1{,}00$ \emph{(base)}      & 73,27 & 659,3 & 479,7 \\
  \bottomrule
\end{tabular}
\end{center}

Duplicar $h_{\mathrm{ext}}$ de 10 a 20 cuesta 3,6 puntos de rendimiento; llevarlo
de 20 a 30, otros 3,4. Es una sensibilidad apreciable, pero sobre un parámetro
cuyo rango de incertidumbre es justamente ese. La emisividad tiene un efecto
parecido y en el mismo sentido: más emisividad, más pérdidas y pared más fría.

\section{Límites conocidos}

\subsection{El intervalo de $T_w$ falla con la placa fría}
\label{sec:fallo}

El extremo inferior de la bisección interna es $T_{\mathrm{cielo}}$, y eso
supone que la pared acaba por encima del sumidero radiante. Es cierto siempre
que la placa esté recibiendo flujo, pero no lo es cuando el flujo es casi nulo y
el aire entra frío: entonces la pared se estabilizaría \emph{por debajo} de
$T_{\mathrm{cielo}}$, fuera del intervalo, y \texttt{\_biseccion} lanza
\texttt{ValueError}.

La condición exacta es que el balance en el extremo inferior sea negativo:
\[
  q_{\mathrm{abs}} \;<\; h_{\mathrm{ext}}\,(T_{\mathrm{cielo}} - T_{\mathrm{amb}})\,A
  \;+\; \frac{A\,(T_{\mathrm{cielo}} - T_f)}{R}.
\]
El primer término es pequeño --- 49~W/m$^2$ con los valores por defecto ---, así que
lo que manda es el segundo: el intervalo falla en cuanto el fluido está más frío
que el cielo y el flujo no basta para compensarlo. Con aire a 17~$^\circ$C y las
constantes por defecto, el umbral está en unos 6,2~kW/m$^2$.

En la práctica esto no afecta a ningún punto de diseño: un receptor de torre
trabaja entre 100 y 500~kW/m$^2$, dos órdenes de magnitud por encima. Sí aparece
al explorar el modelo con flujos residuales y aire ambiente, por ejemplo en un
arranque en frío. Arreglarlo es una línea: bajar el extremo inferior a
\[
  T_w^{\min} = \min\bigl(T_{\mathrm{cielo}},\ T_{\mathrm{amb}},\ T_f\bigr) - 1,
\]
que es la temperatura por debajo de la cual la pared gana calor de todo lo que
la rodea, y el balance es positivo con seguridad.

Conviene señalar que el caso límite que ya se comprobó en su día --- flujo casi
nulo con el aire entrando a 500~$^\circ$C --- sí funciona, y por eso no se había
detectado: con el fluido muy por encima del cielo, el extremo inferior sigue
siendo válido.

\subsection{Comentario desactualizado en el bloque \texttt{\_\_main\_\_}}

El comentario que encabeza el caso de ejemplo dice:

\begin{quote}\small
\emph{Con 300 kW/m2 de pico y este mismo gasto, la pared se va a 946 C y el acero
no aguanta: hay que subir el gasto a 0.6 kg/s.}
\end{quote}

Se escribió cuando el ejemplo usaba $G = 0{,}35$~kg/s, y ahora usa 0,6. Con los
valores actuales la frase se contradice: «este mismo gasto» ya \emph{es} 0,6. Los
números reales son

\begin{center}
\begin{tabular}{llr}
  \toprule
  $q''_{\mathrm{pico}}$ & $G$ & $T_w^{\max}$ \\
  \midrule
  300~kW/m$^2$ & 0,35~kg/s & 951~$^\circ$C \quad(supera el límite) \\
  300~kW/m$^2$ & 0,60~kg/s & 812~$^\circ$C \\
  200~kW/m$^2$ & 0,35~kg/s & 777~$^\circ$C \\
  200~kW/m$^2$ & 0,60~kg/s & 659~$^\circ$C \quad(caso de ejemplo) \\
  \bottomrule
\end{tabular}
\end{center}

El fondo del comentario sigue siendo correcto --- con 0,35~kg/s y 300~kW/m$^2$ el
acero no aguanta, y subir a 0,6 lo arregla ---, solo hay que reescribir la frase
para que no diga «este mismo gasto».

\subsection{Hipótesis físicas}

\begin{itemize}
  \item \textbf{Columnas adiabáticas entre sí}: ni conducción lateral por la
        chapa ni mezcla. Es lo que hace que el perfil de salida reproduzca la
        forma de la campana; con conducción lateral saldría más plano, de modo
        que éste es el caso conservador para el punto caliente.
  \item \textbf{No se calcula la pérdida de carga.} A 47~m/s en un canal de
        pillow plate no va a ser pequeña, y es lo primero que habría que añadir:
        las ecuaciones (18) y (19) ya están en el \texttt{formulas.py} del otro
        paquete, y \texttt{asignacion} ya devuelve $n_1$ y $n_2$, que ahora mismo
        se calculan y no se usan.
  \item \textbf{Reparto de caudal uniforme entre columnas.} En realidad las
        columnas centrales, más calientes y por tanto menos densas, tenderían a
        recibir menos caudal, lo que agrava el punto caliente.
  \item \textbf{Sin reradiación entre porciones} ni factor de vista distinto de
        1: cada porción ve el cielo por completo. En un receptor de cavidad esto
        sería inaceptable; en una placa plana es razonable.
  \item \textbf{La correlación de Nusselt se extrapola en $Pr$.} Se aplica a
        $Re = 3{,}6\cdot10^{4}$, del orden del de los casos con los que se validó,
        pero con aire y no con agua. El $Pr$ del aire, 0,71, queda fuera del
        rango de los ensayos de Piper. Es la extrapolación más discutible del
        modelo y conviene decirlo en la memoria.
  \item \textbf{Régimen estacionario} y sin inercia térmica: no hay transitorios
        de nube ni de arranque.
\end{itemize}

\section{Resumen de uso}

\begin{verbatim}
    import mapa as mp
    import props_fluido as pf
    import receptor as rc

    receptor = rc.Receptor(
        mapa=mp.MapaGaussiano.desde_pico(L=1.5, W=1.0, q_pico=200e3,
                                         sigma_x=1/3, sigma_y=0.5),
        G=0.6,                        # kg/s de aire
        T_ent=pf.CERO_CELSIUS + 300,  # 300 C, uniforme
        p=10e5,                       # 10 bar
        N=20, M=10,
    )
    resultado = receptor.resolver()
    resultado.resumen()
    resultado.dibujar()
\end{verbatim}

El caso de ejemplo del bloque \texttt{\_\_main\_\_} da:

\begin{center}
\begin{tabular}{lr}
  \toprule
  Potencia incidente sobre la placa   & 157,2~kW \\
  Calor al fluido                     & 115,2~kW \\
  Pérdidas                            & 42,0~kW \\
  Rendimiento                         & 73,3~\% \\
  Interceptación del mapa             & 75,1~\% \\
  Salida media / máxima / mínima      & 480 / 546 / 396~$^\circ$C \\
  Pared máxima                        & 659~$^\circ$C \\
  Velocidad en el canal               & 47,5~m/s \\
  Reynolds                            & $3{,}64\cdot10^{4}$ \\
  $h$ interno a la $T$ media          & 728~W/(m$^2$K) \\
  \bottomrule
\end{tabular}
\end{center}

La presión de 10~bar no es un capricho: a presión atmosférica, ese mismo gasto
pediría unos 300~m/s en el canal interno, que no es un punto de trabajo
razonable.

\end{document}
